\begin{lemma}
For $t\geq2$ and sufficiently large prime powers $q$, the polarity graph
$G(t,q)$ has order $n$, degree $d$, and nontrivial spectral parameter
$\lambda$ satisfying
\[
 q^t/2\leq n\leq2q^t,\qquad
 q^{t-1}/2\leq d\leq2q^{t-1},\qquad
 \lambda\leq2\sqrt d.
\]
Its orthogonality relation has no two-row triangular configuration of length
$t+2$: there are no $(a_i,b_i)_{i=1}^{t+2}$ with
$a_i b_i\notin E(G)$ and $a_i b_j\in E(G)$ for all $i<j$.  Moreover, the
old pair digraph $\mathcal D(G)$ has at least $q^{2t}/2$ vertices.
\end{lemma}
\begin{proof}
The points of $G(t,q)$ are the one-dimensional subspaces of
$\mathbb F_q^{t+1}$, so their number is
$(q^{t+1}-1)/(q-1)$, which lies between $q^t/2$ and $2q^t$ for large $q$.
The nonzero solutions of one nontrivial orthogonality equation form a
codimension-one subspace.  After quotienting by nonzero scalar multiples,
this gives degree $(q^t-1)/(q-1)$, between $q^{t-1}/2$ and $2q^{t-1}$.
The common-neighbour calculation for two distinct points gives the usual
identity for the square of the adjacency matrix, so every nonprincipal
eigenvalue has absolute value at most $2\sqrt d$ once $q$ is large.

For the triangular assertion, choose nonzero representatives $x_i,y_i$.
If a nontrivial relation among the $y_i$ has first nonzero coefficient at
$i$, taking its inner product with $x_i$ leaves only the nonzero diagonal
term $\langle x_i,y_i\rangle$, a contradiction.  Thus $t+2$ vectors would
be independent in the $(t+1)$-dimensional ambient space.  Finally, the
ordered nonedge-pair count in the definition of
\noderef{OldPairDigraph}, together with the preceding point and degree
counts, is at least $q^{2t}/2$ for large $q$.
\end{proof}
